\documentclass{article}
\usepackage[letterpaper, textwidth=5.5in, textheight=9.0in]{geometry}
\usepackage{graphicx} % Required for inserting images

% text and typesetting macros
\usepackage{setspace}
\setstretch{1.08}
\renewcommand{\paragraph}[1]{\medskip\par\noindent{\textbf{#1}}~---}
\sloppy\sloppypar\raggedbottom
\frenchspacing
\newcommand{\documentname}{\textsl{Note}}
\newcommand{\sectionname}{Section}
\newcommand{\secref}[1]{\sectionname~\ref{#1}}

% math macros
\usepackage{amsmath}

% collaboration / writing / commenting macros
\usepackage{xcolor}
\newcommand{\hoggsay}[1]{{\color{blue} (Hogg says: #1)}}

\title{\bfseries A custom periodogram for \textsl{Gaia} astrometric residuals.}
\author{Hogg for Ramirez~Delgado and friends}
\date{October 2026}
\begin{document}
\maketitle

\paragraph{Abstract}
We provide a replacement for the periodogram that knows about ``scan angle'' and also knows that the signals we care about are elliptically polarized.

\section{Introduction}
ESA \textsl{Gaia}'s data is fundamentally one-dimensional, and it is entirely in the time domain:
To a very good approximation, \textsl{Gaia} scans the sky and measures the \emph{times} at which stars transit a fiducial point in each of two focal planes.
In this brief \documentname, we ask:
\begin{itemize}
    \item If \textsl{Gaia} does a good job of fitting every object with some kind of kinematic model (parallax and proper motion, say, or even a more sophisticated model), and
    \item if \textsl{Gaia} releases the timing residuals away from that model (or code to compute them), then
    \item can we find periodic signals (of the kinds relevant to astronomy) with some kind of modified periodogram?
\end{itemize}
We are writing this note, so obviously the answer is ``yes.''

\section{Assumptions and choices}
\paragraph{Residual data}
For the purposes of this \documentname, we will assume that a very good kinematic model of every star has been subtracted, and we are left with good residuals.
The model could be generalized, of course.

\paragraph{Generative model}

\paragraph{Noise and robustness}

\section{The method}

\end{document}
